\subsection{The human-approved core and its documentation}
\label{app:verified-census}
\label{app:retained-core-membership}

The human-approved core contains 147 dossiers drawn from 137 distinct sources, selected from BugBank by the rule of Section~\ref{app:core-selection}; human review is not a condition of membership in BugBank itself. The core contains 147 certificate paths and 141 distinct certificate hashes. After the separate checks of reference-line annotations, it supplies 100 primary benchmark inputs and 131 scorable inputs in total; the 31 additional inputs are shorter source texts. Core membership, review evidence, and benchmark eligibility are recorded for each dossier.

Each dossier connects the source text and quoted claim to a concrete witness, its mathematical explanation, and a Lean certificate. Hashes identify the source packet, witness fields, Lean source, and selected declarations; the saved builds confirm that the selected declarations compile, a separate check lists the axioms used by each selected theorem and confirms that only the allowed foundational axioms occur, and the build receipts identify the compiler environment, dependencies, commands, and logs. Replacement witnesses appear alongside their predecessors, so a correction and its reason are inspectable; Section~\ref{app:source-admissibility} develops two such examples.

\subsubsection{Definitions for describing the core}
\label{app:structure-census}

\emph{Source} identifies where the claim appeared. A \emph{mathematical topic} describes the subject of that claim. A \emph{witness object} describes the counterexample itself, such as a graph, number, function, or algorithm trace. These classifications answer different questions: a graph, for example, can expose an error in an algorithms argument.

An \emph{error mechanism} describes the immediate failed inference: a definition or domain condition, a quantifier, an omitted case, an algebraic or order argument, a count, a structural or analytic inference, or an algorithm invariant. The existing mechanism annotations select the operation closest to the quoted error and retain uncertainty flags.

\begin{table}[htbp]
\centering\small
\setlength{\tabcolsep}{3pt}
\begin{minipage}[t]{.48\linewidth}\centering
\begin{tabular}{@{}lrr@{}}
\toprule
Source collection & Dossiers & Inputs \\
\midrule
QEDBench & 25 & 20 \\
Stack Exchange collection & 18 & 11 \\
math.SE collection & 16 & 10 \\
ProofWiki & 28 & 26 \\
arXiv & 15 & 13 \\
Reddit & 42 & 20 \\
Reference-proof additions & 3 & 0 \\
\midrule
Total & 147 & 100 \\
\bottomrule
\end{tabular}
\end{minipage}\hfill
\begin{minipage}[t]{.50\linewidth}\centering
\begin{tabular}{@{}lrr@{}}
\toprule
Error mechanism & Dossiers & Inputs \\
\midrule
Structural inference & 38 & 28 \\
Algebra or order & 37 & 23 \\
Omitted case & 28 & 18 \\
Counting & 14 & 8 \\
Quantifier scope & 11 & 8 \\
Definition or domain & 9 & 6 \\
Algorithm invariant & 8 & 8 \\
Analytic inference & 2 & 1 \\
\midrule
Total & 147 & 100 \\
\bottomrule
\end{tabular}
\end{minipage}
\caption{Composition of the 147-dossier human-approved core and 100 primary
benchmark inputs. Each panel counts every record once. Source collections
follow the release groups, with the two Reddit groups combined. Error
mechanisms use the saved primary annotations; six dossiers, including
three primary inputs, retain a mechanism-annotation uncertainty flag.}
\label{tab:corpus-sources}
\end{table}

\paragraph{Reproduction.}
The repository provides core membership, its inclusion rule, and version evidence in \path{bugbank/core/}. From the repository root, \texttt{python -B bugbank/core/verify.py --release .} reconstructs the core from the saved reviews and checks its scientific content, certificates, and benchmark masks, and \path{benchmark/core/replay.py} reproduces the benchmark scores; the \href{https://github.com/qqliu/CoPEER/blob/main/bugbank/core/README.md}{core README} gives the full commands. Earlier census tables, complete annotation evidence, and coding checks are preserved in \path{bugbank/annotations/}, with portable inspection instructions in its \href{https://github.com/qqliu/CoPEER/blob/main/bugbank/annotations/README.md}{README}.
